\documentclass[11pt]{article}

\usepackage[a4paper,margin=1in]{geometry}
\usepackage{amsmath,amssymb,amsthm,mathtools}
\usepackage[hidelinks]{hyperref}
\hypersetup{
  pdfauthor={Bangyen Pham},
  pdftitle={Displaced-Zero Certificates and the Coefficient Mass of Polynomial Multiples},
  pdfsubject={Lower bounds for the coefficient order statistics and logarithmic coefficient mass of polynomial multiples},
  pdfkeywords={polynomial multiples, coefficient mass, reduced height, Chebyshev systems, exponential sums},
  pdflang={en-US}
}

\newtheorem{theorem}{Theorem}[section]
\newtheorem{lemma}[theorem]{Lemma}
\newtheorem{corollary}[theorem]{Corollary}
\newtheorem{proposition}[theorem]{Proposition}
\newcommand{\R}{\mathbb R}
\newcommand{\N}{\mathbb N}
\DeclareMathOperator{\Tail}{Tail}
\DeclareMathOperator{\sgn}{sgn}

\title{Displaced-Zero Certificates and the Coefficient Mass of Polynomial Multiples}
\author{\shortstack{Bangyen Pham\\
\small Independent Researcher\\
\small\href{mailto:bangyenp@gmail.com}{\texttt{bangyenp@gmail.com}}\\
\small\href{https://orcid.org/0009-0006-9928-2088}{ORCID: 0009-0006-9928-2088}}}
\date{September 2026}

\begin{document}
\maketitle

\begin{abstract}
Fix real roots $2\le r_1\le\cdots\le r_L$, repetitions allowed.  For every
nonzero real or complex polynomial $F=\sum_jf_jx^j$ of degree $D$ divisible by
$\prod_{i=1}^L(x-r_i)$, the $k$-th largest nonleading coefficient
magnitude $b_k$ is at least $|f_D|\prod_{i=k}^L(r_i-1)$ for $1\le k\le L$,
uniformly in $D$.  Summing these rows bounds the logarithmic coefficient
mass $\Lambda(F)=\sum_j\log^+|f_j|$ of a monic multiple below by
$\sum_ii\log(r_i-1)$, which is sharp in order at the first $L$ primes.
With a Jensen-formula bound of order $L^2$, this gives
$\inf\Lambda\asymp L^2+\sum_ii\log(r_i-1)$ over real monic multiples at every
root set with $L\ge2$.  The row $k=2$ is sharp at $(2,3)$, and every monic
multiple of $(x-2)^L$ has $b_k\ge L-k+1$ for $1\le k\le L$.  The key tool is a
tail theorem for normalized exponential sums in the reciprocal roots with
prescribed positive-integer zeros: deleting a displaced largest zero
together with the largest node strictly increases the $\ell^1$ tail, which
yields a product bound over the initial run of consecutive zeros.  As a
by-product, the reduced-height bound $H(P)\ge|P(1)|$ of Dubickas and
Jankauskas holds unconditionally for real roots greater than one,
repetitions allowed.
\end{abstract}

\noindent\textbf{Keywords.} Polynomial multiples; coefficient mass; reduced
height; Chebyshev systems; exponential sums.

\noindent\textbf{2020 Mathematics Subject Classification.} Primary 11C08; Secondary 26C05, 41A50.

\section{Introduction}

For a nonzero $F(x)=\sum_jf_jx^j$ of degree $D$, define its logarithmic
coefficient mass by
\[
  \Lambda(F)=\sum_{j:f_j\ne0}\log^+|f_j|,
  \qquad \log^+x=\max(0,\log x),
\]
the logarithm of the multiplicative height $\prod_{j:f_j\ne0}\max(1,|f_j|)$.
The height $\max_j|f_j|$ sees only one coefficient, and the length
$\sum_j|f_j|$ can be spread over many comparable terms, so neither norm
forces many coefficients to be large at once.  The mass does: it is additive
over coefficients, so lower bounds on several coefficient order statistics
add up to a lower bound for $\Lambda$ that no bound on a single height can
give.  Since $\log^+$
ignores coefficients of magnitude at most one, $\Lambda(\varepsilon F)\to0$
as $\varepsilon\to0$, so every bound below is stated relative to the
leading coefficient $f_D$; the case of interest is the monic one, or more
generally $|f_D|\ge1$, which is also the normalization of the reduced
height~\cite{dj} and costs nothing for integer multiples.

\paragraph{Main results.}
Write $b_k(F)$ for the $k$-th largest of the $D$ nonleading coefficient
magnitudes $|f_j|$, $j<D$, of $F$, counted with repetition, zeros included.
Let $2\le r_1\le\cdots\le r_L$ be real, repetitions allowed, let
$P=\prod_{i=1}^L(x-r_i)$, and let $F$ be a nonzero real multiple of $P$.
Theorem~\ref{thm:order} and
Corollary~\ref{cor:mass} give
\[
  b_k(F)\ge|f_D|\prod_{i=k}^{L}(r_i-1)
  \quad(1\le k\le L),
  \qquad
  \Lambda(F)\ge(L+1)\log|f_D|+\sum_{i=1}^{L}i\log(r_i-1).
\]
We call the first inequality at a given $k$ the row $k$.  Both bounds are
uniform in the degree and in any further roots of the multiple, and they
extend to complex multiples (Corollaries~\ref{cor:complex} and~\ref{cor:mass}).  The row $k=1$
holds for all $r_1>1$.  For the first $L$ primes and $|f_D|\ge1$ the mass bound is
$\Lambda(F)\ge(\frac12+o(1))L^2\log L$ (Corollary~\ref{cor:prime}), and the
infimum over monic multiples is
$\frac12L^2\log L+\frac12L^2\log\log L+O(L^2)$ (Corollary~\ref{cor:infimum}).
For every root set the product $P$ itself comes within $O(L^2)$ of the
infimum of $\Lambda$ over monic multiples (Corollary~\ref{cor:near}).  As
the roots approach $2$ the mass bound degenerates; a separate argument by
Jensen's formula gives $\Lambda(F)\ge L^2/80$ for monic multiples when
$L\ge8$ (Proposition~\ref{prop:quadratic}), and together
\[
  \inf_{Q\text{ monic}}\Lambda(PQ)\ \asymp\ L^2+\sum_{i=1}^Li\log(r_i-1)
  \qquad(L\ge2),
\]
with absolute constants (Corollary~\ref{cor:allroots}).

\paragraph{Comparison with other bounds.}
The Mahler measure gives a larger constant in the first row, but it loses
a factor $\sqrt D$ and bounds only one coefficient.  A monic $F$ has
$M(F)\ge\prod_ir_i$, and Landau's inequality $M(F)\le\|F\|_2$ gives
$b_1(F)\ge\bigl((\prod_ir_i)^2-1\bigr)^{1/2}D^{-1/2}$, which gains the
factor $\prod_i\frac{r_i}{r_i-1}$ over the row $k=1$.  For the mass it gives only
$\Lambda(F)\ge\sum_i\log r_i-\frac12\log(D+1)$, at most of order $L\log L$
at the first $L$ primes and vacuous once $\log D\gg L\log L$.  Jensen's
formula on a disk of radius $R<1$, the tool behind the zero counts of
Borwein, Erd\'elyi and K\'os for polynomials with bounded
coefficients~\cite{bek}, removes the degree: exempting the $k-1$ largest
magnitudes gives the bound~\eqref{eq:jensen} of
Proposition~\ref{prop:quadratic}, which uses only $r_i\ge2$.  The two
methods are complementary.  Near the root $2$ the certificate rows tend to
$1$, while~\eqref{eq:jensen} gives $b_1\ge2^L/(eL)-1/(L-1)$ and
exponentially large first rows.  The right side of~\eqref{eq:jensen} is
below $2^{L-k+1}$ for every $R$, however, and below $1/L$ at $k=L$.  The
certificate row $k$ is therefore the stronger bound whenever
$\prod_{i=k}^L(r_i-1)\ge2^{L-k+1}$, in particular whenever $r_k\ge3$; at the
first $L$ primes it is stronger in every row, and
Corollaries~\ref{cor:prime}, \ref{cor:infimum} and~\ref{cor:near} rest on
it.  In the last rows every certificate row is at least $1$, and at the
root $2$ Theorem~\ref{thm:everyrow} raises the row $k$ to $L-k+1$.

\paragraph{Method.}
To bound $b_k$, we exclude any $k-1$ positions and build a normalized
exponential sum in the reciprocal roots, the certificate, vanishing there
and at enough of the top positions $D-1,D-2,\ldots$
(Lemma~\ref{lem:slack}).  For monic $F$, with the excluded positions
carrying $b_1,\ldots,b_{k-1}$, pairing the coefficients of $F$ with the
certificate $w$ at the roots gives $1\le b_k(F)\Tail(w)$, where $\Tail(w)$
is the $\ell^1$ tail of $w$ (see~\eqref{eq:dual}).  This tail is
bounded, uniformly over zero
sets with a given leading run, by Theorem~\ref{thm:tail}, whose
engine is a strict deletion step: deleting a displaced largest zero together
with the largest node, the reciprocal of the smallest root, strictly
increases the tail.  Section~\ref{sec:sharp} uses two further tools:
Jensen's formula for the first rows near the root $2$, and confluent limits
of the certificates at the single root $2$ of multiplicity $L$.

\paragraph{Relation to earlier work.}
Lemma~\ref{lem:consecutive} and Theorem~\ref{thm:tail} remove the
same-sign hypothesis under which Dubickas and Jankauskas~\cite{dj} evaluate
their dual bound for the reduced height, and extend their consecutive-zero
evaluation to a tail bound for arbitrary integer zero sets, which gives the
rows $k\ge2$.  Reduced length and reduced height minimize the two classical
coefficient
norms over normalized multiples~\cite{dubickas,schinzel,schinzel-ii,schinzel-complex,dj};
the Euclidean analogue is effectively computable~\cite{frw}.  Integer,
bounded-height and sparse multiples, and power series with restricted
coefficients and a root of prescribed multiplicity, form related
lines~\cite{mignotte,bombieri-vaaler,bbbp-multiple,bbbp-ray,grt}.  Closest to
this paper, Dubickas and Jankauskas~\cite{dj} associate to reciprocal roots
and a prescribed zero set $J$ a normalized determinant sequence $S_n(J)$,
which is the exponential sum $w$ below; the reciprocal of its tail is their
dual lower bound for the reduced height, their inequality~(8).  They
evaluate the consecutive case $J=\{1,\ldots,c-1\}$ on the condition that
all $S_n$ with $n$ large have the same sign~\cite[p.~333]{dj}, a hypothesis
kept through their Corollary~7~\cite[Section~2]{dj} and checked case by
case~\cite[Sections~3--4]{dj}.  Here $H(P)$ is the reduced height, the infimum of the
heights of the multiples $PQ$ with $Q$ monic.  Lemma~\ref{lem:consecutive}
proves the same-sign condition for all distinct nodes in $(0,1)$, so their
bound $H(P)\ge|P(1)|$, which follows from the row $k=1$ of Theorem~\ref{thm:order},
holds unconditionally for real roots greater than one, repetitions
allowed.

\paragraph{Sharpness and companion papers.}
Section~\ref{sec:sharp} asks how close the rows are to optimal.  The loss
in the row $k=1$ is unbounded: if $r_1\ge2$, every monic
multiple has $b_1\ge(r_1+\cdots+r_L)/2$, which exceeds $|P(1)|$ by a factor of order $L$
at suitable distinct roots (Proposition~\ref{prop:unbounded}), and the
Jensen bound of Proposition~\ref{prop:quadratic} gives
$b_1\ge2^L/(eL)-1/(L-1)$ once $L\ge3$, so near the root $2$ the loss is
exponential in $L$.  The row
$k=2$ is an infimum at $(2,3)$ (Proposition~\ref{prop:sharp23}), so no constant
factor improves the rows $k\ge2$ as a family; the companion
paper~\cite{coefficient-mass-attainment} characterizes the root sets at which
a row is an infimum by a partial-sum criterion, repeated roots included.
Every monic multiple of $(x-2)^L$ has $b_k\ge L-k+1$ for every row
$1\le k\le L$ (Theorem~\ref{thm:everyrow}).  Its proof removes the degree
from the problem (Section~\ref{sec:dfree}), and the companion
paper~\cite{coefficient-mass-rows} builds on that reduction to determine the
rows at a single root $r\ge2$ exactly: the row $k$ equals the top row of a
multiple of $(x-r)^{L-k+1}$, which sharpens Theorem~\ref{thm:everyrow} at
$r=2$.  All roots here are real; the
companion paper~\cite{coefficient-mass-complex} treats complex roots, where
full orbits of $x\mapsto\zeta x$, in particular purely imaginary pairs,
reduce exactly to real roots, and $K$ roots of modulus at least $R$ in a
sector of half-width $\delta$ cost a mass of order
$\min(K^2,K/\delta)\log R$ for multiples of degree $O(K)$.

\paragraph{Formalization.}
Sections~2 and~3 through Corollary~\ref{cor:mass}, that is,
Theorem~\ref{thm:order} and Corollary~\ref{cor:mass} with the whole chain
leading to them from the zero bound, are formally verified in Lean~4 with
Mathlib, including the equality case of Theorem~\ref{thm:tail} and the
relaxation of $r_1\ge2$ to $r_1>1$ in the first row of Lemma~\ref{lem:slack}
and Theorem~\ref{thm:order}; the worked examples are not.  Of Section~\ref{sec:sharp}, Proposition~\ref{prop:sharp23} and
Lemma~\ref{lem:rowcert} through Theorem~\ref{thm:everyrow} are
formalized; Propositions~\ref{prop:unbounded}
and~\ref{prop:quadratic} and Corollary~\ref{cor:allroots} are not.  The formal proofs, in the
directory \texttt{CoefficientMass} of the repository
\href{https://github.com/bangyen/coefficient-mass}{bangyen/coefficient-mass},
follow the ones given here, and the formal proof of
Corollary~\ref{cor:mass} depends only on the standard axioms of Lean
(propositional extensionality, choice and quotient soundness).  The zero
bound is formalized for distinct zeros, the sign changes at prescribed
zeros being derived there from weak alternation.

\paragraph{Conventions.}
All logarithms are natural, $\N_{>0}=\{1,2,\ldots\}$, and every asymptotic
statement is for $L\to\infty$ with absolute implied constants unless
parameters are named.

\section{A displaced-zero tail theorem}

Let $\rho_1>\cdots>\rho_c\ge2$ and $y_i=\rho_i^{-1}$, so that
$0<y_1<y_2<\cdots<y_c\le1/2$.  For any $Z\subset\N_{>0}$ of size $c-1$, the
generalized Vandermonde system
\[
  w_d=\sum_{i=1}^c a_i y_i^d,
  \qquad w_0=1,
  \qquad w_z=0\quad(z\in Z),
\]
determines the coefficients $a_i$ uniquely.  Indeed, a kernel vector would be
a nonzero $c$-node exponential sum with the $c$ zeros $\{0\}\cup Z$,
contradicting the zero bound below.  It is convenient to regard the sequence
$w_0,w_1,\ldots$ as the restriction to $d\in\{0,1,2,\ldots\}$ of the
real-analytic function $d\mapsto\sum_ia_iy_i^d$; all statements about zeros
and signs of an exponential sum refer to this function of the real variable.

\paragraph{Zero bound.}
A nonzero exponential sum on $n$ distinct positive nodes has at most $n-1$
real zeros, counted with multiplicity: write it as $\sum_ia_ie^{\lambda_id}$
with distinct $\lambda_i$, factor out one exponential, and induct on $n$ by
Rolle's theorem, the derivative being a sum on one fewer
node~\cite[Part~Five, Ch.~1]{polya-szego}; equivalently, the $e^{\lambda d}$
form a Chebyshev system~\cite{karlin}.
Consequently a nonzero $n$-node sum with $n-1$
distinct prescribed real zeros has exactly those zeros, each is simple, and the sum
changes sign at each of them.

A form of the zero bound that needs no local analysis at the zeros is
\emph{weak alternation}: if a nonzero sum $f$ on $r$ distinct positive nodes
satisfies $\varepsilon(-1)^kf(t_k)\ge0$ at points $t_0<t_1<\cdots<t_n$, for a
fixed sign $\varepsilon$, then $n+1\le r$.  Divide by the exponential of one
node, giving a function $g$ that keeps the hypothesis, and apply the mean
value theorem to $g$ on each $[t_{k-1},t_k]$: the derivative, a sum on the other $r-1$ nodes, satisfies
$\varepsilon(-1)^kg'(s_k)\ge0$ at some $s_k\in(t_{k-1},t_k)$, so it weakly
alternates at the $n$ points $s_1<\cdots<s_n$, and induction on $r$ applies.
If the derivative vanishes identically, $f$ is a single exponential, which has
a strict sign, so $n=0$.  A zero at which $f$ keeps its sign is a weak
alternation on both sides, so this bound does the work of a count with
multiplicity.

\begin{lemma}[Consecutive zero set]\label{lem:consecutive}
Let $0<y_1<\cdots<y_c<1$; no cutoff $y_c\le1/2$ is imposed.  Let $w$ be the
exponential sum on these nodes with $w_0=1$ and $w_z=0$ for $z\in\{1,\ldots,c-1\}$;
it exists and is unique because a kernel vector of the interpolation system
would be a nonzero $c$-node sum with the $c$ zeros $\{0,\ldots,c-1\}$,
contradicting the zero bound, which needs only distinct positive nodes.  Then
\[
  w_d=(-1)^{c+1}\Bigl(\prod_iy_i\Bigr)h_{d-c}(y_1,\ldots,y_c)
  \qquad(d\ge c),
\]
where $h_j$ is the complete homogeneous symmetric polynomial, so these terms
are one-signed and
\[
  \Tail(w):=\sum_{d\ge1}|w_d|=\prod_{i=1}^c\frac{y_i}{1-y_i}.
\]
\end{lemma}

\begin{proof}
Put
\[
  U(x)=\sum_{d\ge0}w_dx^d=\frac{N(x)}{\prod_i(1-y_ix)},
  \qquad \deg N\le c-1.
\]
The first $c-1$ nonconstant coefficients vanish, so $N(x)-\prod_i(1-y_ix)$
is a constant multiple of $x^c$; that constant is fixed by the vanishing
$x^c$ coefficient of $N$ (which has degree at most $c-1$), giving
\[
  N(x)=\prod_i(1-y_ix)+(-1)^{c+1}\Bigl(\prod_i y_i\Bigr)x^c .
\]
Since $1/\prod_i(1-y_ix)=\sum_{j\ge0}h_j(y_1,\ldots,y_c)x^j$, the displayed
formula for $w_d$ follows, and summing the one-signed terms gives
\[
  \Tail(w)=\Bigl(\prod_i y_i\Bigr)\sum_{j\ge0}h_j(y_1,\ldots,y_c)
          =\prod_i\frac{y_i}{1-y_i}.
\]
\end{proof}

\begin{theorem}[Displaced-zero tail bound]\label{thm:tail}
Let $\rho_1>\cdots>\rho_c\ge2$, $y_i=\rho_i^{-1}$, $Z\subset\N_{>0}$ with
$|Z|=c-1$, and let $w_d=\sum_{i=1}^ca_iy_i^d$ be the exponential sum
with $w_0=1$ and $w_z=0$ for $z\in Z$.  Put
$\ell=\max\{0\le k\le c-1:\{1,\ldots,k\}\subseteq Z\}$, the leading run of
$Z$; the elements of $Z$ outside $\{1,\ldots,\ell\}$ are its displaced zeros.
Then
\begin{equation}\label{eq:tail}
  \Tail(w)=\sum_{d\ge1}|w_d|
  \le \prod_{i=1}^{\ell+1}\frac{y_i}{1-y_i}
  =\frac{1}{\prod_{i=1}^{\ell+1}(\rho_i-1)},
\end{equation}
a product over the $\ell+1$ largest $\rho_i$.  Equality holds exactly when $Z=\{1,\ldots,c-1\}$.
\end{theorem}

The proof removes the displaced zeros one at a time, starting with the
largest, each time together with the largest node $y_c=1/\rho_c$, and shows
that each removal strictly increases the tail.  It does not compare zero
sets at fixed nodes, and no comparison of that kind holds in general.  At the nodes
$(\rho_1,\rho_2)=(3,2)$ the zero set $Z=\{1\}$
gives $w_d=3^{1-d}-2^{1-d}$ and $\Tail(w)=\frac12$, equality in
\eqref{eq:tail}, while $Z=\{2\}$ gives $w_d=\frac95\,3^{-d}-\frac45\,2^{-d}$
and $\Tail(w)=\frac3{10}$; at $(\rho_1,\rho_2)=(5,3)$, however, $Z=\{1\}$
gives $\frac18$ and $Z=\{2\}$ gives $\frac9{64}$.

The hard step is the following bound: for the
correction sum $G$, the ratio $E_d/\widehat V_d$ with $E_d=G_d-y_cG_{d-1}$
never exceeds its value at $d=m+1$, just past the marked points; it is
proved by weak alternation.

\begin{lemma}[Ratio bound for the correction sum]\label{lem:interlace}
Let $c\ge2$ and $0<y_1<\cdots<y_c$, let $Z'\subset\N_{>0}$ satisfy $|Z'|=c-2$, and put
$m=\max(Z'\cup\{0\})$; the points of $\{0\}\cup Z'$ are the marked
points.  Let $\widehat V$ be the exponential sum on
$y_1,\ldots,y_{c-1}$ vanishing on $Z'$, and $G$ the sum on
$y_1,\ldots,y_c$ vanishing on $\{0\}\cup Z'$.  Each is determined by those
conditions up to a scalar and, by the zero bound, has constant sign for
$d>m$; normalize both to be positive there.  Write $E_d=G_d-y_cG_{d-1}$,
$\sigma=G_{m+1}/\widehat V_{m+1}$, and
\[
  \Delta_d=\sigma\widehat V_d-E_d.
\]
Then $\Delta_d\ge0$ for every $d\ge m+1$.  No upper bound on the nodes is
needed.
\end{lemma}

\begin{proof}
The sums $E$ and $\widehat V$ both live on the $c-1$ nodes
$y_1,\ldots,y_{c-1}$: writing $G_d=\sum_ia_iy_i^d$ gives
\[
  E_d=\sum_{i=1}^{c-1}a_i(y_i-y_c)y_i^{\,d-1}
      =\sum_{i=1}^{c-1}\beta_iy_i^d,
\]
with $\beta_i=a_i(1-y_c/y_i)$, the $i=c$ term having cancelled.  Hence $\Delta$
is an exponential sum on the $c-1$ nodes $y_1,\ldots,y_{c-1}$, some of its
coefficients possibly zero.

We record its values at the points of $Z'$ and at $m+1$.  First,
$\Delta_{m+1}=0$, because $G_m=0$ gives $E_{m+1}=G_{m+1}=\sigma\widehat
V_{m+1}$.  Next let $p\in Z'$.  Then $\widehat V_p=G_p=0$, so
$\Delta_p=-E_p=y_cG_{p-1}$.  By the zero bound $G$ has exactly the simple
zeros $\{0\}\cup Z'$ and is positive past $m$, so on the component of
$(0,\infty)\setminus Z'$ just below $p$ it has the sign $(-1)^{a(p)}$, where
$a(p)=|\{\alpha\in Z':\alpha\ge p\}|$.  Hence $\Delta_p$ is zero, when $p-1$
is marked, or has the sign $(-1)^{a(p)}$.

Suppose $\Delta_d<0$ for some $d>m+1$.  List the $c$ points
$Z'\cup\{m+1,d\}$ in increasing order and let $k(p)$ be the position of $p$,
counted from $0$; thus $k(p)=|\{\alpha\in Z':\alpha<p\}|=c-2-a(p)$ for
$p\in Z'$, $k(m+1)=c-2$ and $k(d)=c-1$.  Then
$(-1)^c(-1)^{k(p)}\Delta_p\ge0$ at every listed point: at $p\in Z'$ by the
sign of $\Delta_p$, at $m+1$ because $\Delta_{m+1}=0$, and at $d$ because
$(-1)^c(-1)^{c-1}\Delta_d=-\Delta_d>0$.  So $\Delta$, which is not
identically zero since $\Delta_d\ne0$, weakly alternates at $c$ points
while having at most $c-1$ nodes, contradicting weak alternation.  Hence
$\Delta_d\ge0$ for every $d\ge m+1$.
\end{proof}

\begin{proof}[Proof of Theorem~\ref{thm:tail}]
Induct on the number $c-1-\ell$ of displaced zeros.

\emph{Base case.}  For $Z=\{1,\ldots,c-1\}$ this is
Lemma~\ref{lem:consecutive}; the cutoff $y_c\le1/2$ is not used there.

\emph{Displaced step.}  Suppose $\ell<c-1$ and set $z=\max Z$, so
$z\ge c\ge\ell+2$.  Put $Z'=Z\setminus\{z\}$ and $m=\max(Z'\cup\{0\})$.  Let
$V$ be the sum on $y_1,\ldots,y_{c-1}$ with $V_0=1$ and $V|_{Z'}=0$, and $G$
the sum on all $c$ nodes with $G|_{\{0\}\cup Z'}=0$ and $G_z=1$.  Then
$w=V-V_zG$, both sides satisfying the same $c$ interpolation conditions.
The set $Z'$ has the same leading run $\ell$ as $Z$, so by induction
$\Tail(V)$ is bounded by the right side of \eqref{eq:tail}, and it remains
to prove $\Tail(w)<\Tail(V)$.

\emph{Sign comparison.}  By the zero bound, $G$ has exactly the simple zeros
$\{0\}\cup Z'$, $V$ exactly $Z'$, and $w$ exactly $Z$.  Write
$\eta=\sgn(V_z)$ and $\mu=|V_z|$, so $w=V-\eta\mu G$.  Since $V$ and $G$ flip
at the same points, $\sgn(V_d/G_d)$ is constant on $(0,\infty)\setminus Z'$,
equal to its value $\eta$ at $d=z$; hence
$w_d=\eta\sgn(G_d)\bigl(|V_d|-\mu|G_d|\bigr)$ for $d>0$.  The sign of
$\varphi(d)=|V_d|-\mu|G_d|$ is thus that of $\eta\sgn(G_d)w_d$, which does
not change across $Z'$, where $w$ and $G$ both flip, and does change at $z$,
where only $w$ flips.  As $\varphi(0)=1$ and $\varphi$ has no zero in
$(0,\min Z)$, we get
\[
  |V_d|>\mu|G_d|\quad(1\le d<z,\ d\notin Z'),
  \qquad
  |V_d|<\mu|G_d|\quad(d>z),
\]
with $|w_d|=\bigl||V_d|-\mu|G_d|\bigr|$ throughout and equality at the
zeros of $w$.

For $X\in\{V,G\}$ write
\[
  \Gamma_X(z)=\frac{\sum_{d\ge z}|X_d|}{|X_z|},
  \qquad \Sigma_G=\sum_{1\le d<z}|G_d|.
\]
Splitting $\Tail(V)-\Tail(w)$ at $z$ and using the two displayed
inequalities, $|V_z|=\mu|G_z|$ and $|G_z|=1$, gives
\begin{equation}\label{eq:split}
  \Tail(V)-\Tail(w)
  =\mu\bigl(2\Gamma_V(z)-\Gamma_G(z)+\Sigma_G\bigr).
\end{equation}

\emph{The correction ratio.}  We claim
\begin{equation}\label{eq:gamma}
  \Gamma_G(z)\le\frac{\Gamma_V(z)}{1-y_c}\le2\Gamma_V(z).
\end{equation}
Let $\widehat V=\eta V$, which is positive at every $d>m$, and
$E_d=G_d-y_cG_{d-1}$, an exponential sum on the first $c-1$ nodes.  We first
show
\begin{equation}\label{eq:drive}
  \frac{E_d}{\widehat V_d}\le
  \frac{G_z}{\widehat V_z}\qquad(d\ge z>m).
\end{equation}

Set $\sigma=G_{m+1}/\widehat V_{m+1}$ and
$\Delta_d=\sigma\widehat V_d-E_d$.  Both $\widehat V$ and $G$ are positive for
$d>m$, the latter because $G_z=1$ with $z>m$, so Lemma~\ref{lem:interlace}
applies and gives $\Delta_d\ge0$ for $d\ge m+1$, that is,
$E_d/\widehat V_d\le\sigma$ for $d\ge m+1$.  Next, $\sigma\le G_z/\widehat V_z$, with equality
when $z=m+1$.  So let $z>m+1$, put $\tau=G_z/\widehat V_z>0$, and consider
the $c$-node sum
\[
  \Theta_d=G_d-\tau\widehat V_d .
\]
It vanishes on $Z'$ and at $z$, which is $c-1$ points, and
$\Theta_0=-\tau\widehat V_0=-\tau\eta\ne0$, so by the zero bound $\Theta$
changes sign exactly at those points.  The $c-2$ points of $Z'$ lie in
$(0,m+1)$ and $z$ does not, so $\sgn\Theta_{m+1}=(-1)^{c}\sgn\Theta_0
=-(-1)^{c}\eta$.  And $\eta=(-1)^{c}$, since $V_0=1$ and $V$ flips at each of
the $c-2$ points of $Z'$, all below $z$.  Hence $\Theta_{m+1}<0$, that is,
$\sigma<\tau$.  Therefore
\[
  \frac{E_d}{\widehat V_d}\le\sigma\le\frac{G_z}{\widehat V_z}
  \qquad(d\ge z>m),
\]
which is \eqref{eq:drive}.

Induction in $G_d=y_cG_{d-1}+E_d$, with $G_z=1$ and
\eqref{eq:drive}, gives for $t\ge0$
\[
  \frac{G_{z+t}}{G_z}
  \le\sum_{j=0}^t y_c^{t-j}
       \frac{\widehat V_{z+j}}{\widehat V_z}.
\]
Summing over $t$ (a geometric series in $y_c$) proves the first inequality
in \eqref{eq:gamma}; the second uses $y_c\le1/2$.

\emph{Strictness.}  It remains to see that $\Sigma_G>0$.  If $z>m+1$, then
$G_{m+1}\ne0$.  If $z=m+1$, then $|Z'|=c-2<z-1$, so some $1\le d<z$ avoids
$\{0\}\cup Z'$ and $G_d\ne0$.  So \eqref{eq:split} is positive and
$\Tail(w)<\Tail(V)$: deleting the largest prescribed zero together with the
largest node $y_c=1/\rho_c$ strictly increases the tail.  This completes the induction, and
equality in \eqref{eq:tail} occurs only in the base case.
\end{proof}

The restriction $\rho_i\ge2$ in Theorem~\ref{thm:tail} matters: for $y=(\frac1{10},\frac45)$ and
$Z=\{2\}$ one finds $w_d=\frac{64}{63}10^{-d}-\frac1{63}(\frac45)^d$ and
$\Tail(w)=\frac{52}{405}>\frac19$, so \eqref{eq:tail} fails.  The cutoff
enters only through $1/(1-y_c)\le2$ in \eqref{eq:gamma} and the factors
discarded in the proof of Lemma~\ref{lem:slack}.

\section{Coefficient order statistics}

Let
\[
 P(x)=\prod_{i=1}^L(x-r_i),\qquad 2\le r_1<\cdots<r_L,
\]
with the roots distinct for the lemma below (repetitions are handled in
Theorem~\ref{thm:order}), and let $F(x)=\sum_{j=0}^Df_jx^j\in\R[x]$ be a
nonzero multiple of $P$.  As in the introduction, $b_1\ge b_2\ge\cdots$ are
the magnitudes $|f_j|$, $j<D$, in decreasing order.

\begin{lemma}[Certificate with excluded positions]\label{lem:slack}
Assume the roots $2\le r_1<\cdots<r_L$ are distinct and $|f_D|=1$.  Fix $0\le u\le L-1$, with $r_1>1$
in place of $r_1\ge2$ permitted when $u=0$.  For every set $U$ of $u$
coefficient positions below $D$, some $j<D$ outside $U$ satisfies
\[
 |f_j|\ge\prod_{i=u+1}^L(r_i-1).
\]
\end{lemma}

\begin{proof}
Use all $c=L$ roots $\rho_1=r_L>\rho_2=r_{L-1}>\cdots>\rho_c=r_1$.  A
position $j$ has backward distance $d=D-j$.  Let $W\subseteq\N_{>0}$ be the
set of distances of the positions in $U$, and let $w$ be the certificate on
these $c$ nodes, with $w_0=1$ and with zero set $Z$ consisting of $W$
together with the first $L-u-1$ positive distances outside $W$.  The $u$
positions have distinct distances, so $|W|=u$, and $L-u-1\ge0$ because
$u\le L-1$; hence $|Z|=c-1$, as required; write $w_d=\sum_i a_iy_i^d$.

The dual-certificate identity is explicit.  Since $F(\rho_i)=0$ for the $c$
chosen roots,
\[
  0=\sum_{i=1}^c a_i\rho_i^{-D}F(\rho_i)
   =\sum_{j\le D}f_jw_{D-j}
   =f_D+\sum_{j<D}f_jw_{D-j},
\]
using $w_0=1$.  The terms with $j\in U$ vanish because $D-j\in W\subseteq Z$.
Hence, as $|f_D|=1$,
\begin{equation}\label{eq:dual}
  1\le\sum_{j<D,\ j\notin U}|f_j|\,|w_{D-j}|
  \le\Bigl(\max_{j<D,\ j\notin U}|f_j|\Bigr)
        \sum_{j<D,\ j\notin U}|w_{D-j}|.
\end{equation}

The set $Z$ contains $\{1,\ldots,L-u-1\}$, since such a distance lies in $W$
or among the first $L-u-1$ distances outside $W$, so its leading run is at least
$L-u-1$, and Theorem~\ref{thm:tail} bounds $\Tail(w)$ by a product of at
least $L-u$ factors $y_i/(1-y_i)$ over the smallest reciprocal nodes.  Each
factor is at most $1$ because $\rho_i\ge2$, so we may keep only the first
$L-u$, which belong to $\rho_1,\ldots,\rho_{L-u}$, that is, to
$r_L,\ldots,r_{u+1}$.  For $u=0$ nothing is discarded, $Z=\{1,\ldots,L-1\}$,
and Lemma~\ref{lem:consecutive} needs only $r_1>1$.  Therefore
\[
  \sum_{j<D,\ j\notin U}|w_{D-j}|
  \le\Tail(w)
  \le\prod_{i=1}^{L-u}\frac{\rho_i^{-1}}{1-\rho_i^{-1}}
  =\frac{1}{\prod_{i=u+1}^L(r_i-1)},
\]
and \eqref{eq:dual} gives the claim.
\end{proof}

\begin{theorem}[Coefficient order statistics]\label{thm:order}
Let $r_1\le\cdots\le r_L$ be real with $r_1\ge2$, repetitions allowed, and
let $F$ be a nonzero real multiple of $\prod_{i=1}^L(x-r_i)$ of degree $D$,
and let $b_1\ge b_2\ge\cdots\ge b_D$ be the magnitudes $|f_j|$, $j<D$, in
decreasing order.  Then for every $0\le u\le L-1$,
\begin{equation}\label{eq:order}
  b_{u+1}\ge|f_D|\prod_{i=u+1}^L(r_i-1).
\end{equation}
If instead only $r_1>1$ is assumed, \eqref{eq:order} still holds for $u=0$.
\end{theorem}

\begin{proof}
Assume first that the $r_i$ are distinct and $|f_D|=1$, and put
$T=\prod_{i=u+1}^L(r_i-1)$.  Take for $U$ the positions of $b_1,\ldots,b_u$.
Lemma~\ref{lem:slack} gives some $j<D$ outside $U$ with $|f_j|\ge T$, and
every position outside $U$ has magnitude at most $b_{u+1}$, so
$b_{u+1}\ge T$.

For general $f_D\ne0$, the quotient $F/f_D$ has leading coefficient $1$ and
coefficient magnitudes $b_k(F)/|f_D|$, so the normalized case gives
$b_{u+1}(F)\ge|f_D|\,T$.

For repetitions, put $r_{i,\varepsilon}=r_i+i\varepsilon$ with
$\varepsilon>0$, so that $2\le r_{1,\varepsilon}<\cdots<r_{L,\varepsilon}$ and
$r_{i,\varepsilon}\to r_i$, and write $F=Q\prod_{i=1}^L(x-r_i)$ and
$F_\varepsilon=Q\prod_{i=1}^L(x-r_{i,\varepsilon})$ with the same $Q$.  Then
$\deg F_\varepsilon=D$, the leading coefficient is still $f_D$, and each
coefficient of $F_\varepsilon$ is a polynomial in $\varepsilon$, so the
coefficient vector of $F_\varepsilon$ tends to that of $F$; an order statistic of a converging tuple is continuous, so
$b_{u+1}(F_\varepsilon)\to b_{u+1}(F)$.  The distinct case gives
$b_{u+1}(F_\varepsilon)\ge|f_D|\prod_{i=u+1}^L(r_{i,\varepsilon}-1)$, and letting
$\varepsilon\to0$ proves \eqref{eq:order}.

The $u=0$ clause follows from the $u=0$ clause of Lemma~\ref{lem:slack} by
the same perturbation, which keeps $r_{i,\varepsilon}>1$.
\end{proof}

\begin{corollary}[Complex multiples]\label{cor:complex}
Theorem~\ref{thm:order} holds verbatim for a nonzero $F\in\mathbb C[x]$
divisible by $\prod_{i=1}^L(x-r_i)$.
\end{corollary}

\begin{proof}
Write $g=F/f_D=PQ$ with $Q\in\mathbb C[x]$ monic.  Since $P$ is real,
$\operatorname{Re}g=P\operatorname{Re}Q$ is a monic real multiple of $P$ of
the same degree $D$, and $|\operatorname{Re}g_j|\le|g_j|$ for every $j$.
Coefficientwise domination passes to order statistics, so
Theorem~\ref{thm:order} applied to $\operatorname{Re}g$ gives
$b_k(g)\ge b_k(\operatorname{Re}g)\ge\prod_{i=k}^L(r_i-1)$; multiplying
through by $|f_D|$ gives \eqref{eq:order} for $F$.
\end{proof}

\begin{corollary}[Logarithmic mass]\label{cor:mass}
Let $r_1\le\cdots\le r_L$ be real with $r_1\ge2$, repetitions allowed, and
let $F$ be a nonzero real or complex multiple of $\prod_{i=1}^L(x-r_i)$.  Then
\begin{equation}\label{eq:mass}
 \Lambda(F)\ge(L+1)\log|f_D|+\sum_{i=1}^{L}i\log(r_i-1).
\end{equation}
\end{corollary}

\begin{proof}
By Theorem~\ref{thm:order}, with Corollary~\ref{cor:complex} for complex
$F$, \eqref{eq:order} holds.  Each product in it is at least one because $r_i\ge2$, so the
$L$ largest nonleading magnitudes are nonzero, as is $f_D$; since $D\ge L$
these are $L+1$ distinct positions.  Summing $\log^+\ge\log$ over them alone,
\[
 \Lambda(F)\ge\log|f_D|+\sum_{u=0}^{L-1}\log b_{u+1}
 \ge(L+1)\log|f_D|+\sum_{i=1}^{L}i\log(r_i-1),
\]
using $\log b_{u+1}\ge\log|f_D|+\sum_{i=u+1}^L\log(r_i-1)$ and collecting the
index-$i$ terms, which occur exactly $i$ times.
\end{proof}

At the bottom of the allowed range the two bounds behave differently.  For
$r_1=\cdots=r_L=2$ every product in \eqref{eq:order} equals $1$, so
Theorem~\ref{thm:order} still forces $L$ nonleading coefficients of
magnitude at least $|f_D|$, while \eqref{eq:mass} reduces to
$\Lambda(F)\ge(L+1)\log|f_D|$ and says nothing more.  The order-statistic
bound stays informative over the whole root range; the mass bound does not.
Section~\ref{sec:sharp} shows that this degeneration is in the bound and not
in $\Lambda$: for monic multiples at $r_i=2$,
Proposition~\ref{prop:quadratic} gives $b_1\ge2^L/(eL)-1/(L-1)$ and a
mass of order $L^2$, and Theorem~\ref{thm:everyrow} the last rows
$b_k\ge L-k+1$.

\begin{corollary}[Prime roots]\label{cor:prime}
Let $r_1\le\cdots\le r_L$ with $r_i\ge p_i$, the $i$th prime, and let $F$ be
a nonzero real multiple of $\prod_{i=1}^L(x-r_i)$ with $|f_D|\ge1$.  Then
\[
 \Lambda(F)\ge\Bigl(\frac12+o(1)\Bigr)L^2\log L .
\]
In particular $\Lambda(F)=\Omega(L^2\log L)$, and the bound is asymptotically
attained: for $F=P_L=\prod_{i=1}^L(x-p_i)$,
$\Lambda(P_L)\le(\tfrac12+o(1))L^2\log L$.
\end{corollary}

\begin{proof}
Drop the nonnegative term $(L+1)\log|f_D|$ from \eqref{eq:mass}.  For
$i\ge2$ we have $r_i-1\ge p_i-1\ge i$, so
$\Lambda(F)\ge\sum_{i=2}^Li\log i=(\frac12+o(1))L^2\log L$.  For the upper
bound, the nonleading coefficients of $P_L$ are the elementary symmetric
sums $e_k(p_1,\ldots,p_L)\le\binom{L}{k}p_L^k$, so
\[
 \Lambda(P_L)\le\sum_{k=1}^L\log\binom{L}{k}
 +\frac{L(L+1)}2\log p_L
 \le L^2\log2+\frac{L(L+1)}2\log p_L,
\]
which is $(\tfrac12+o(1))L^2\log L$ because $\log p_L\sim\log L$.
\end{proof}

\begin{corollary}[Sharp prime-root infimum]\label{cor:infimum}
Let $P_L=\prod_{i=1}^L(x-p_i)$ and
\[
  M_L=\inf_{\substack{Q\in\R[x]\\ Q\text{ monic}}}\Lambda(QP_L).
\]
Then
\[
  M_L=\frac12L^2\log L+\frac12L^2\log\log L+O(L^2).
\]
\end{corollary}

\begin{proof}
Corollary~\ref{cor:mass} with $r_i=p_i$ and $|f_D|=1$ gives, for every monic
multiple $F$ of $P_L$,
\[
  \Lambda(F)\ge\sum_{2\le i\le L}i\log(p_i-1)
  =\frac12L^2\log L+\frac12L^2\log\log L+O(L^2),
\]
the $i=1$ term being zero.  The equality holds because
$i\log i<p_i<i(\log i+\log\log i)$ for $i\ge6$~\cite{rosser-schoenfeld}
gives $\log(p_i-1)=\log i+\log\log i+O(1)$ for $i\ge2$, while
$\sum_{2\le i\le L}i\log i=\frac12L^2\log L+O(L^2)$ and
$\sum_{2\le i\le L}i\log\log i=\frac12L^2\log\log L+o(L^2)$.  The matching upper
bound is the one for $P_L$ from Corollary~\ref{cor:prime},
\[
  \Lambda(P_L)\le L^2\log2+\frac{L(L+1)}2\log p_L
  =\frac12L^2\log L+\frac12L^2\log\log L+O(L^2).
\]
\end{proof}

The prime roots are not special in this respect: at every root set the
product $P$ itself comes within $O(L^2)$ of the infimum, so
Corollary~\ref{cor:mass} determines the infimum up to $O(L^2)$.

\begin{corollary}[Near-optimality of the root product]\label{cor:near}
Let $2\le r_1\le\cdots\le r_L$, repetitions allowed,
$P=\prod_{i=1}^L(x-r_i)$, and
\[
  M(P)=\inf_{\substack{Q\in\R[x]\\ Q\text{ monic}}}\Lambda(PQ).
\]
Then
\begin{equation}\label{eq:near}
  0\le\Lambda(P)-M(P)
  \le\sum_{k=1}^L\log\binom Lk
    +\sum_{i=1}^L i\log\frac{r_i}{r_i-1}
  =O(L^2),
\end{equation}
with an absolute implied constant.
\end{corollary}

\begin{proof}
Taking $Q=1$ gives $M(P)\le\Lambda(P)$, while Corollary~\ref{cor:mass} gives
$M(P)\ge\sum_{i=1}^L i\log(r_i-1)$, because a monic multiple has $|f_D|=1$.
In the other direction the nonleading coefficients of $P$ are the elementary
symmetric sums, and each $k$-fold product of roots is at most the product of
the $k$ largest, so
\[
  e_k(r_1,\ldots,r_L)\le\binom Lk\prod_{i=L-k+1}^{L}r_i .
\]
Since $e_k\ge\binom Lk2^k>1$, summing logarithms gives
\[
  \Lambda(P)=\sum_{k=1}^L\log e_k(r_1,\ldots,r_L)
  \le\sum_{k=1}^L\log\binom Lk+\sum_{i=1}^L i\log r_i ,
\]
the sum over $k$ collecting $\log r_i$ exactly $i$ times.
Subtracting the lower bound and using
$\log r_i-\log(r_i-1)=\log\frac{r_i}{r_i-1}$ gives \eqref{eq:near}; finally
$\binom Lk\le2^L$ and $r_i/(r_i-1)\le2$ because $r_i\ge2$, so the right side
is at most $(L^2+\frac{L(L+1)}2)\log2=O(L^2)$.
\end{proof}

\section{Sharpness}\label{sec:sharp}

How close is Theorem~\ref{thm:order} to optimal?  We treat the first row,
the bottom of the root range, the rows $k\ge2$, and finally, at the root
$2$, a degree-free form of the rows.

\subsection{The first row}

For a monic multiple, row $k=1$ of \eqref{eq:order} is
$b_1\ge\prod_i(r_i-1)=|P(1)|$, which contains the bound $H(P)\ge|P(1)|$ of
Dubickas and Jankauskas~\cite{dj}.  It can be far from optimal.

\begin{proposition}[Unbounded looseness of the row $k=1$]
\label{prop:unbounded}
Let $1<r_1\le\cdots\le r_L$, repetitions allowed, and put
$B=\prod_{i=1}^L(r_i-1)$ and $S=\sum_{i=1}^Lr_i$.  Then every monic multiple
$F$ of $\prod_i(x-r_i)$ satisfies
\[
  b_1(F)\ \ge\ \frac{SB}{B+1},
\]
whose ratio to the value $B=|P(1)|$ given by \eqref{eq:order} at $u=0$ is
$S/(B+1)$.  If moreover $r_i\ge2$ for every $i$, then $B\ge1$ and hence
\begin{equation}\label{eq:halfsum}
  b_1(F)\ \ge\ \frac S2=\frac{r_1+\cdots+r_L}2 .
\end{equation}
For $r_i=2+iL^{-3}$ the ratio $S/(B+1)$ is $\sim L$, so the ratio
between a certificate bound and $|P(1)|$ is unbounded in $L$.
\end{proposition}

\begin{proof}
Assume first that the roots are distinct.  Take $Z=\{2,\ldots,L\}$, of size
$L-1$, and let $w$ be the certificate on the
nodes $y_i=r_i^{-1}$, in any order, as nothing below depends on it, so $w_0=1$ and $w_2=\cdots=w_L=0$.  Write
$W(t)=\sum_{d\ge0}w_dt^d=N(t)/\Pi(t)$ with $\Pi(t)=\prod_i(1-y_it)$ and
$\deg N\le L-1$.  Then $N(t)-(1+w_1t)\Pi(t)$ has degree at most $L+1$ and is
divisible by $t^{L+1}$, hence equals $\kappa t^{L+1}$, so
\[
  W(t)=1+w_1t+\frac{\kappa t^{L+1}}{\Pi(t)} .
\]
The coefficient of $t^L$ in $N$ vanishes, which reads
\[
  (-1)^Le_L(y)+w_1(-1)^{L-1}e_{L-1}(y)=0,
  \qquad\text{while}\qquad
  e_{L-1}(y)=e_L(y)\sum_{i}y_i^{-1}=e_L(y)\,S,
\]
so that $w_1=1/S$.  Then
$\kappa=-w_1[t^L]\Pi(t)=(-1)^{L+1}e_L(y)/S$, and $\Pi(t)=e_L(y)\prod_i(r_i-t)$,
so the last term is $(-1)^{L+1}t^{L+1}/\bigl(S\prod_i(r_i-t)\bigr)$.
Expanding $1/\Pi(t)=\sum_{j\ge0}h_j(y)t^j$ gives
$w_{L+1+j}=(-1)^{L+1}e_L(y)h_j(y)/S$, which is one-signed, so
\[
  \Tail(w)=w_1+\frac{e_L(y)}S\sum_{j\ge0}h_j(y)
  =\frac1S+\frac1S\prod_{i}\frac{y_i}{1-y_i}
  =\frac{1+B^{-1}}S .
\]
The dual-certificate identity in the proof of Lemma~\ref{lem:slack}, applied
with $U=\varnothing$ and this $w$, gives $b_1(F)\ge1/\Tail(w)=SB/(B+1)$.  The
computation used only $0<y_i<1$, that is $r_i>1$.

Repetitions follow by the perturbation in the proof of
Theorem~\ref{thm:order}, as $B$ and $S$ are continuous in the roots.  If
every $r_i\ge2$ then $B\ge1$, so $B/(B+1)\ge\tfrac12$, which is
\eqref{eq:halfsum}.  Finally, for $r_i=2+iL^{-3}$ one has
$\log B=\sum_i\log(1+iL^{-3})=O(L^{-1})$, so $B\to1$, while
$S=2L+O(L^{-1})$; hence $S/(B+1)\sim L$.
\end{proof}

\subsection{The bottom of the range}\label{sec:bottom}

The infimum of the mass over monic multiples is of order $L^2$ throughout
the range where \eqref{eq:mass} degenerates, by Jensen's formula.  At
$r_i=2$ the certificate of Proposition~\ref{prop:unbounded} gives
$b_1\ge L$, where the row $k=1$ of \eqref{eq:order} gives only $1$, and
this too is far from the truth: the first rows grow exponentially in $L$
(the remark after Corollary~\ref{cor:allroots}).  The proof uses
Jensen's formula rather than a certificate, in the spirit
of~\cite{bbbp-multiple,bbbp-ray} and of the zero counting at $1$ for
polynomials with bounded coefficients in~\cite{bek}, with $k-1$ coefficients exempted rather
than all being bounded.

\begin{proposition}[Quadratic mass at every root set]
\label{prop:quadratic}
Let $2\le r_1\le\cdots\le r_L$, repetitions allowed.  Every monic real
multiple $F$ of $\prod_{i=1}^L(x-r_i)$ satisfies, for every
$\tfrac12<R<1$,
\begin{equation}\label{eq:jensen}
  b_k(F)\ \ge\ \frac{(2R)^{L-k+1}-1}{(1-R)^{-k}-1},\qquad 1\le k\le L,
\end{equation}
and hence $\Lambda(F)\ge L^2/80$ for $L\ge8$.
\end{proposition}

\begin{proof}
Let $D=\deg F$ and let $F^*(t)=t^DF(1/t)=\sum_{d=0}^Df_{D-d}t^d$ be the
reciprocal polynomial, so $F^*(0)=f_D=1$, while
$F^*(1/r_i)=r_i^{-D}F(r_i)=0$, so $F^*$ has at least $L$ zeros in $(0,\tfrac12]$ counted with multiplicity,
namely the $1/r_i$.

For a real polynomial $h$ and an integer $z\ge1$ put $T_zh=h-(t/z)h'$, so that
$[t^d](T_zh)=(1-d/z)[t^d]h$: the operator annihilates the coefficient at
$t^z$ and fixes the constant term.  It also satisfies
\[
  T_zh=-\frac{t^{z+1}}z\bigl(t^{-z}h\bigr)' .
\]
If $h$ has $m$ zeros in an interval $I\subset(0,\infty)$, counted with
multiplicity, then so does $t^{-z}h$, and by Rolle's theorem its derivative
has at least $m-1$ zeros in the convex hull of them, hence in $I$; the factor
$-t^{z+1}/z$ contributes no zero in $I$.  So $T_zh$ retains at least $m-1$
zeros in $I$.

Fix $k$ and let $z_1<\cdots<z_{k-1}$ be the backward distances $D-j$ of the
positions $j<D$ carrying the $k-1$ largest nonleading magnitudes.  Put
$g=T_{z_{k-1}}\cdots T_{z_1}F^*$, so that
\[
  [t^d]g=[t^d]F^*\prod_{i=1}^{k-1}\Bigl(1-\frac d{z_i}\Bigr).
\]
Then $g(0)=1$, the coefficients at $d=z_1,\ldots,z_{k-1}$ vanish, and by the
previous paragraph $g$ has at least $L-k+1$ zeros in $(0,\tfrac12]$.

If $d\ge1$ and $d\notin\{z_i\}$ then $|[t^d]F^*|=|f_{D-d}|\le b_k(F)$, that
position not being among the $k-1$ largest.  The $z_i$ are distinct positive
integers, so $z_i\ge i$ and
\[
  \prod_{i=1}^{k-1}\Bigl|1-\frac d{z_i}\Bigr|
  \le\prod_{i=1}^{k-1}\Bigl(1+\frac di\Bigr)=\binom{d+k-1}{k-1} .
\]
Hence $|[t^d]g|\le b_k(F)\binom{d+k-1}{k-1}$ for $d\ge1$, and since
$\sum_{d\ge0}\binom{d+k-1}{k-1}x^d=(1-x)^{-k}$, taking $x=R$ gives
\[
  \max_{|t|=R}|g(t)|\ \le\ 1+b_k(F)\bigl((1-R)^{-k}-1\bigr).
\]
Jensen's formula for $g$ on $|t|\le R$, with $g(0)=1$ and at least
$L-k+1$ zeros of modulus at most $\tfrac12$, gives
$(2R)^{L-k+1}\le\max_{|t|=R}|g|$, which is \eqref{eq:jensen}.

Now take $R=\tfrac9{10}$, let $L\ge8$ and $1\le k\le\lfloor L/8\rfloor$,
and put $\lambda=L/8\ge1$.  Then $L-k+1\ge7\lambda\ge7$ and $10^k-1<10^\lambda$.
Since $9^7=4782969>61\cdot5^7$, we have $(9/5)^7>61$, so
$(9/5)^{L-k+1}-1\ge(9/5)^{7\lambda}\bigl(1-\tfrac1{61}\bigr)$ and
\[
  b_k(F)\ \ge\ \frac{60}{61}\Bigl(\frac{(9/5)^7}{10}\Bigr)^{\lambda}
  \ \ge\ \frac{60}{61}\Bigl(\frac{61}{50}\Bigr)^{\lambda}5^\lambda\ \ge\ 5^\lambda\ >\ e^{L/5},
\]
the last step because $e<\tfrac{68}{25}$ and $68^8<5^5\cdot25^8$ give $e^8<5^5$.  These $\lfloor L/8\rfloor$
magnitudes sit at distinct positions and exceed $1$, and
$\lfloor L/8\rfloor\ge L/16$, so $\Lambda(F)\ge(L/16)(L/5)=L^2/80$.
\end{proof}

Combining this with the certificate bound closes the mass problem to within
absolute constants over the whole root range.

\begin{corollary}[The mass infimum at every root set]\label{cor:allroots}
Let $L\ge2$, let $2\le r_1\le\cdots\le r_L$, put $P=\prod_{i=1}^L(x-r_i)$ and
$M(P)=\inf\{\Lambda(PQ):Q\in\R[x]\text{ monic}\}$.  Then
\[
  M(P)\ \asymp\ L^2+\sum_{i=1}^Li\log(r_i-1),
\]
with absolute implied constants.
\end{corollary}

\begin{proof}
For the lower bound, Proposition~\ref{prop:quadratic} gives
$M(P)\ge L^2/80$ when $L\ge8$, and \eqref{eq:halfsum} gives
$b_1\ge L\ge2$, hence $M(P)\ge\log2\ge L^2\log2/49$, when $2\le L\le7$.
Corollary~\ref{cor:mass} with $|f_D|=1$ gives
$M(P)\ge\sum_ii\log(r_i-1)$, so $M(P)$ is at least a constant multiple of
the sum of the two terms.
For the upper bound take $Q=1$: as in Corollary~\ref{cor:near},
$\Lambda(P)\le\sum_k\log\binom Lk+\sum_ii\log r_i$, and
$\sum_ii\log r_i\le\sum_ii\log(r_i-1)+\tfrac{L(L+1)}2\log2$ because
$r_i/(r_i-1)\le2$, so
$\Lambda(P)\le2L^2\log2+\sum_ii\log(r_i-1)$.
\end{proof}

At $k=1$ and $R=1-1/L$, $L\ge3$, \eqref{eq:jensen} reads
$b_1\ge((2-2/L)^L-1)/(L-1)$, and $(1-1/L)^{L-1}\ge e^{-1}$ gives
\[
  b_1(F)\ \ge\ \frac{2^L}{eL}-\frac1{L-1},
\]
exponentially more than the value $1$ of \eqref{eq:order} and the value $L$
of Proposition~\ref{prop:unbounded} at $r_i=2$, a growth the
certificates of Lemma~\ref{lem:slack} do not see.  The bound
\eqref{eq:jensen} is useful only while $(1-R)^{-k}$ has not overtaken
$(2R)^{L-k+1}$, that is, for $k$ up to a constant multiple of $L$; it says
nothing about the last rows.

For the last rows at the root $2$, Theorem~\ref{thm:everyrow} at the end of
this section proves every row $b_k\ge L-k+1$ for monic multiples of
$(x-2)^L$.  This is weaker than \eqref{eq:jensen} for the first rows
(numerically, for $k\le3$ at $L=20$ and $k\le7$ at $L=40$), and its
content is the last rows, where \eqref{eq:jensen} is empty.

\subsection{The rows \texorpdfstring{$k\ge2$}{k at least 2}}

For $u\ge1$ the excluded set $U$ in Lemma~\ref{lem:slack} is arbitrary, and
Lemma~\ref{lem:slack} places the remaining $L-u-1$ zeros of the certificate
at the first available distances; that choice is convenient, not optimal.
Still, the rows are infima at some root sets.

\begin{proposition}[Sharpness of the row $k=2$ at $(2,3)$]
\label{prop:sharp23}
For $n\ge2$ put
\[
  t_n=\frac{2}{1-2^{1-n}+3^{-n}},
  \qquad a_n=2+t_n\bigl(1-2^{-n}\bigr),
  \qquad
  F_n(x)=x^{n+1}-a_nx^n+t_n\sum_{j=0}^{n-1}x^j .
\]
Then each $F_n$ is a monic multiple of $(x-2)(x-3)$ with $b_2(F_n)=t_n$, and
$t_n$ decreases to $2$.  Hence
\[
  \inf_{\substack{Q\in\R[x]\\ Q\text{ monic}}}
  b_2\bigl(Q(x)(x-2)(x-3)\bigr)=2,
\]
the value \eqref{eq:order} asserts for $u=1$.  Every $F_n$ has $b_2>2$.
\end{proposition}

\begin{proof}
Summing the geometric progression, $F_n(2)=0$ and $F_n(3)=0$ read
$a_n=2+t_n(1-2^{-n})$ and $a_n=3+\tfrac12t_n(1-3^{-n})$; the two agree exactly
when $t_n\bigl(\tfrac12-2^{-n}+\tfrac123^{-n}\bigr)=1$, which is the stated
$t_n$.  Writing $d_n=1-2^{1-n}+3^{-n}$, one has $d_n<1$ because
$3^{-n}<2\cdot2^{-n}$, and $d_n\ge\tfrac12$ because $2^{1-n}\le\tfrac12$ for
$n\ge2$, so $2<t_n\le4$; and
$d_{n+1}-d_n=2^{-n}-\tfrac23\cdot3^{-n}>0$, so $d_n$ increases to $1$ and
$t_n$ decreases to $2$.  Finally $a_n-t_n=2-t_n2^{-n}\ge2-4\cdot2^{-n}>0$ for
$n\ge2$, so the nonleading magnitudes are $a_n$ once and $t_n$ exactly $n$
times with $a_n>t_n>0$; thus $b_2(F_n)=t_n$.  Theorem~\ref{thm:order} with
$(r_1,r_2)=(2,3)$ and $u=1$ gives $b_2\ge2$ for every monic multiple, and
$t_n\to2$.
\end{proof}

At $(\tfrac52,3,\tfrac72)$ the row $k=2$ is not an infimum, so the rows
are not infima at every root set.  The companion
paper~\cite{coefficient-mass-attainment} characterizes the root sets at
which \eqref{eq:order} is an infimum: for $2\le r_1\le\cdots\le r_L$,
repeated roots included, the row $k=u+1$ is the infimum of $b_{u+1}(PQ)$
over monic $Q$ exactly when every partial sum of the coefficients of
$P_u=\prod_{i>u}(x-r_i)$, read down from the leading one, has absolute value
at most $|P_u(1)|$.  At $(\tfrac52,3,\tfrac72)$ with $u=1$ the
partial sum $1-\tfrac{13}2=-\tfrac{11}2$ of the two leading coefficients of
$P_1=(x-3)(x-\tfrac72)$ exceeds $P_1(1)=5$ in absolute value.

\subsection{A \texorpdfstring{$D$}{D}-free form of the rows}\label{sec:dfree}

We return to the rows $b_k\ge L-k+1$ for monic multiples of
$(x-2)^L$ announced at the end of Section~\ref{sec:bottom}.  At a
single root of multiplicity $L$ the certificates of Section~3 are replaced
by their confluent limits, polynomials in $s$ times $2^{-s}$, and the
degree $D$ drops out of the problem.

For a finite set $Z\subset\N_{>0}$ put
\begin{equation}\label{eq:phi}
  r_Z(s)=\prod_{z\in Z}\Bigl(1-\frac sz\Bigr),\qquad
  \Phi(Z)=\sum_{s\ge1}|r_Z(s)|\,2^{-s},
\end{equation}
so $r_Z(0)=1$, $\deg r_Z=|Z|$, and $\Phi(\varnothing)=1$.  The rows reduce to
the following statement about $\Phi$.

\begin{lemma}[Certificate for a row]\label{lem:rowcert}
Let $F$ be a monic multiple of $(x-2)^L$ of degree $D$, let $1\le k\le L$,
let $S$ be the set of backward distances $D-j$ of $k-1$ positions $j<D$
carrying the $k-1$ largest nonleading magnitudes, and put $n=L-k+1$.  If
some finite $Z\supseteq S$ with $|Z|\le L-1$ has $\Phi(Z)\le1/n$, with
$\Phi$ as in \eqref{eq:phi}, then $b_k(F)\ge n$.
\end{lemma}

\begin{proof}
Write $\theta=x\,d/dx$ and $\psi=r_Z$, of degree at most $L-1$.  Since
$\psi(D-\theta)x^j=\psi(D-j)x^j$ and each application of $\theta$ lowers the order
of the zero of $F$ at $2$ by at most one, $(\psi(D-\theta)F)(2)=0$, that is
$\sum_{j\le D}f_j\psi(D-j)2^{j}=0$.  Dividing by $2^D$ and using $f_D=\psi(0)=1$
and $\psi=0$ on $S$,
\[
  1=\Bigl|\sum_{1\le s\le D,\ s\notin S}f_{D-s}\,\psi(s)2^{-s}\Bigr|
   \le b_k(F)\sum_{s\ge1}|\psi(s)|2^{-s}=b_k(F)\,\Phi(Z)\le\frac{b_k(F)}n,
\]
because every position outside the $k-1$ largest has magnitude at most
$b_k(F)$.
\end{proof}

For $n\ge1$ consider the statement
\begin{equation}\tag{$\ast_n$}\label{eq:starn}
  \text{for every finite $A\subset\N_{>0}$ there is a finite $Z\supseteq A$
  with $|Z|\le|A|+n-1$ and $\Phi(Z)\le1/n$.}
\end{equation}
By Lemma~\ref{lem:rowcert} with $A=S$, row $k$ at $L$ follows, for every
$D$, from \eqref{eq:starn} with $n=L-k+1$.  Two tools evaluate $\Phi$.

\begin{lemma}[Deleting the largest zero]\label{lem:confdel}
Let $\Phi$ be as in \eqref{eq:phi}.  If $Z\ne\varnothing$ and
$Z\ne\{1,\ldots,|Z|\}$, then $\Phi(Z)\le\Phi(Z\setminus\{\max Z\})$.  Moreover
$\Phi(\{1,\ldots,\ell\})=1$ for every $\ell\ge0$, so $\Phi(Z)\le1$ for every
$Z$.
\end{lemma}

\begin{proof}
Put $c=|Z|+1$ and $Z'=Z\setminus\{\max Z\}$.  For small $\varepsilon>0$ take
the nodes $y_i=\tfrac12-(c-i)\varepsilon$, so $0<y_1<\cdots<y_c=\tfrac12$,
let $w^\varepsilon$ be the sum on $y_1,\ldots,y_c$ with $w_0=1$ and
$w|_Z=0$, and $V^\varepsilon$ the sum on $y_1,\ldots,y_{c-1}$ with $V_0=1$
and $V|_{Z'}=0$.  As $Z$ is not $\{1,\ldots,c-1\}$, the deletion step in the
proof of Theorem~\ref{thm:tail} gives $\Tail(w^\varepsilon)<\Tail(V^\varepsilon)$.

Now let $\varepsilon\to0$.  Write $w^\varepsilon_d=\sum_{i=1}^c\alpha_i
[y_1,\ldots,y_i](y^d)$ in Newton form, the bracket being the divided
difference of $y\mapsto y^d$, equal to $h_{d-i+1}(y_1,\ldots,y_i)$.  The
$c$ interpolation conditions at $d\in\{0\}\cup Z$ form a linear system for
$\alpha$ whose matrix tends to the matrix with entries
$\binom d{i-1}2^{-(d-i+1)}$; that limit is invertible, because
$\binom d{i-1}$, $i\le c$, is a basis of the polynomials of degree below
$c$ and such a polynomial vanishing at the $c$ points $\{0\}\cup Z$ is zero.
Hence $\alpha$ stays bounded and converges, and since
$0\le h_{d-i+1}(y_1,\ldots,y_i)\le\binom d{i-1}2^{-(d-i+1)}$, the terms
$|w^\varepsilon_d|$ are dominated by $Cd^{\,c}2^{-d}$.  The limit
$w^0_d=p(d)2^{-d}$ has $\deg p<c$, $p(0)=1$ and $p|_Z=0$, so $p=r_Z$, and by
dominated convergence $\Tail(w^\varepsilon)\to\Phi(Z)$.  Likewise
$\Tail(V^\varepsilon)\to\Phi(Z')$, which proves the inequality.

Finally $|r_{\{1,\ldots,\ell\}}(s)|=\binom{s-1}\ell$ for $s>\ell$ and $0$ for
$1\le s\le\ell$, and $\sum_{s>\ell}\binom{s-1}\ell2^{-s}=1$.  Deleting the
largest element repeatedly from any $Z$ stops at its leading run
$\{1,\ldots,\ell\}$, never passing through an initial segment before, so
$\Phi(Z)\le1$.
\end{proof}

The last clause is the row $k=L$ ($n=1$) in $D$-free form.  For the other
rows we compute $\Phi$ exactly on sets that are full up to a few holes.

\begin{lemma}[Hole integral]\label{lem:holeint}
Let $C\subset\N_{>0}$ be finite and nonempty, $m=\max C$, and for
$c\in C$ put $\pi_c=\prod_{c'\in C}c'\big/\prod_{c'\in C\setminus\{c\}}(c-c')$.
For every $N\ge m$, with $\Phi$ as in \eqref{eq:phi},
\[
  \Phi\bigl(\{1,\ldots,N\}\setminus C\bigr)=\int_0^1v^{N-m}M_C(v)\,dv,
  \qquad
  M_C(v)=\sum_{c\in C}2^{-c}v^{m-c}\bigl(\pi_c(1+v)^{c-1}+|\pi_c|(1-v)^{c-1}\bigr).
\]
\end{lemma}

\begin{proof}
Let $Y=\{1,\ldots,N\}\setminus C$ and $\psi=r_Y$.  For $s>N$ every factor of
$\psi(s)$ is negative and
$|\psi(s)|=\binom{s-1}N\prod_{c\in C}\frac c{s-c}=\binom{s-1}N\sum_{c\in C}\frac{\pi_c}{s-c}$,
the last step by partial fractions.  Using $\frac1{s-c}=\int_0^1u^{s-c-1}du$,
$\sum_{s>N}\binom{s-1}N y^s=y^{N+1}(1-y)^{-N-1}$ and then $v=u/(2-u)$,
\[
  \sum_{s>N}\binom{s-1}N\frac{2^{-s}}{s-c}
  =\int_0^1\frac{u^{N-c}\,du}{(2-u)^{N+1}}
  =2^{-c}\int_0^1v^{N-c}(1+v)^{c-1}\,dv .
\]
For $s\le N$, $\psi(s)=0$ unless $s\in C$, and for $c\in C$,
\[
  |\psi(c)|=\frac{\prod_{p\le N,\,p\ne c}|p-c|/p}{\prod_{c'\in C\setminus\{c\}}|c'-c|/c'}
  =\frac1{\binom Nc}\cdot\frac{|\pi_c|}c
  =|\pi_c|\int_0^1v^{N-c}(1-v)^{c-1}\,dv .
\]
Summing the finitely many series gives the formula.
\end{proof}

For $C=\{1\}$ and $C=\{2\}$ one finds $M_C\equiv1$, so
\begin{equation}\label{eq:onehole}
  \Phi(\{2,\ldots,N\})=\frac1N,\qquad
  \Phi(\{1,\ldots,N\}\setminus\{2\})=\frac1{N-1};
\end{equation}
the first is the identity
$\sum_{s\ge1}|\prod_{z=2}^N(1-s/z)|2^{-s}=1/N$ behind
Proposition~\ref{prop:unbounded}.  For $C=\{3,b\}$ with $b\ge4$,
$\pi_b=-\pi_3=3b/(b-3)$, and with $(1+v)^2-(1-v)^2=4v$,
\begin{equation}\label{eq:twohole}
  M_{\{3,b\}}(v)=\frac{3b}{b-3}\Bigl(2^{-b}\bigl((1+v)^{b-1}+(1-v)^{b-1}\bigr)
                   -\tfrac12v^{b-2}\Bigr).
\end{equation}

\begin{lemma}[Two holes]\label{lem:twohole}
For $b\ge4$ and $q\ge0$, with $\Phi$ as in \eqref{eq:phi},
$\Phi(\{1,\ldots,b+q\}\setminus\{3,b\})\le1/(q+1)$.
\end{lemma}

\begin{proof}
For $b=4$, \eqref{eq:twohole} is $\tfrac32(1-v^2)$, and
$\int_0^1v^q\tfrac32(1-v^2)\,dv=\frac3{(q+1)(q+3)}\le\frac1{q+1}$.  Let
$b\ge5$ and expand $M=M_{\{3,b\}}$ in the Bernstein basis
$B_i=\binom{b-1}iv^i(1-v)^{b-1-i}$, $0\le i\le b-1$.  From
$1+v=(1-v)+2v$ we get
$2^{-b}(1+v)^{b-1}=\sum_i2^{-(b-i)}B_i$; also $(1-v)^{b-1}=B_0$ and
$v^{b-2}=B_{b-1}+B_{b-2}/(b-1)$.  So $M=\frac{3b}{b-3}\sum_i\gamma_iB_i$ with
\[
  \gamma_0=2^{1-b},\quad
  \gamma_i=2^{-(b-i)}\ (1\le i\le b-3),\quad
  \gamma_{b-2}=\frac{b-3}{4(b-1)},\quad
  \gamma_{b-1}=0 .
\]
Every coefficient $\frac{3b}{b-3}\gamma_i$ is at most $1$: the largest
candidates are $\frac{3b}{8(b-3)}\le1$ (as $b\ge5$) and $\frac{3b}{4(b-1)}\le1$.
Since the $B_i$ are nonnegative and sum to $1$, $0\le M\le1$ on $[0,1]$, and
Lemma~\ref{lem:holeint} gives
$\Phi\le\int_0^1v^q\,dv=1/(q+1)$.
\end{proof}

\begin{proposition}[Rows away from the top three positions]\label{prop:rowsfree}
Let $n\ge1$ and let $A\subset\N_{>0}$ be finite with
$\{1,2,3\}\not\subseteq A$.  Then some finite $Z\supseteq A$ has
$|Z|\le|A|+n-1$ and $\Phi(Z)\le1/n$.  Consequently, if $F$ is a monic
multiple of $(x-2)^L$ of degree $D$ and $1\le k\le L$, and the $k-1$ largest
nonleading magnitudes can be chosen at positions whose set does not contain all
three of $D-1,D-2,D-3$, then $b_k(F)\ge L-k+1$.
\end{proposition}

\begin{proof}
For $n=1$ take $Z=A$ and use Lemma~\ref{lem:confdel}.  Let $n\ge2$, let $\ell$
be the leading run of $A$ (as in Theorem~\ref{thm:tail}; so $\ell\le2$), and let $c_1=\ell+1<c_2<\cdots$ be the
positive integers outside $A$.  Keep the holes
$K=\{1\}$, $\{2\}$, $\{3,c_2\}$ for $\ell=0,1,2$, fill the next $n-1$ holes
after $K$, and put
\[
  N=c_{|K|+n-1},\qquad Y=\{1,\ldots,N\}\setminus K,\qquad Z=Y\cup A .
\]
Then $Y\supseteq A\cap[1,N]$, the elements of $Z\setminus A$ are the $n-1$
filled holes, and $|Z|\le|A|+n-1$.

\emph{From $Z$ to $Y$.}  List the elements of $Z$ above $N$ increasingly and
add them to $Y$ one at a time.  Every set so obtained has the hole
$\min K\le3$ below an element larger than it, so it is not of the form
$\{1,\ldots,t\}$; reading the additions backwards, Lemma~\ref{lem:confdel}
gives $\Phi(Z)\le\Phi(Y)$.

\emph{The base.}  For $\ell=0$, $Y=\{2,\ldots,N\}$ with $N=c_n\ge n$, and
\eqref{eq:onehole} gives $\Phi(Y)=1/N\le1/n$.  For $\ell=1$,
$Y=\{1,\ldots,N\}\setminus\{2\}$ with $N=c_n\ge n+1$, and
$\Phi(Y)=1/(N-1)\le1/n$.  For $\ell=2$ put $b=c_2\ge4$; the $n-1$ filled holes
exceed $b$, so $N\ge b+n-1$.  The sets
$\{1,\ldots,t\}\setminus\{3,b\}$, $b+n-1\le t\le N$, contain $b+1$ and miss
$3$, so Lemma~\ref{lem:confdel} and Lemma~\ref{lem:twohole} with $q=n-1$ give
$\Phi(Y)\le\Phi(\{1,\ldots,b+n-1\}\setminus\{3,b\})\le1/n$.

The consequence is Lemma~\ref{lem:rowcert} with $A=S$: here $|S|=k-1$, so
$|Z|\le L-1$, and $\{1,2,3\}\not\subseteq S$ is the hypothesis on the
positions.
\end{proof}

\paragraph{The top three positions.}
There remains \eqref{eq:starn} for $A\supseteq\{1,2,3\}$, that is, a leading run
$\ell\ge3$.  The construction is the one above with the first $\ell+1$ holes kept.
The two parts of $\Phi$ that appear in the proof of Lemma~\ref{lem:holeint}
are bounded separately: the values at the kept holes by comparison with the
consecutive block $\{\ell+1,\ldots,2\ell+1\}$, and the tail beyond $N$ by a Beta
integral.  Throughout, $\ell\ge0$, $C=\{c_1<\cdots<c_{\ell+1}\}\subset\N_{>0}$ with
$c_1=\ell+1$, $m=c_{\ell+1}$, $q\ge0$, $N=m+q$, $Y=\{1,\ldots,N\}\setminus C$ and
$\psi=r_Y$.  Then $m\ge2\ell+1$, and by the proof of Lemma~\ref{lem:holeint}
$\Phi(Y)=\Phi_{\mathrm{hole}}(C,q)+\Phi_{\mathrm{tail}}(C,q)$, where
\[
  \Phi_{\mathrm{hole}}(C,q)=\sum_{c\in C}|\psi(c)|\,2^{-c},\qquad
  \Phi_{\mathrm{tail}}(C,q)=\sum_{s>N}|\psi(s)|\,2^{-s}
        =\sum_{s>N}\binom{s-1}N2^{-s}\prod_{c\in C}\frac c{s-c}.
\]

\begin{lemma}[Values at the kept holes]\label{lem:holevals}
Let $\ell\ge0$ and $q\ge0$, let $C=\{c_1<\cdots<c_{\ell+1}\}\subset\N_{>0}$
with $c_1=\ell+1$, and put $m=c_{\ell+1}$, $N=m+q$,
$Y=\{1,\ldots,N\}\setminus C$ and $\psi=r_Y$.  Then
$\Phi_{\mathrm{hole}}(C,q)=\sum_{c\in C}|\psi(c)|\,2^{-c}\le\dfrac1{2(q+1)}$.
\end{lemma}

\begin{proof}
Fix $1\le i\le \ell+1$, put $c=c_i$ and $i'=\ell+1-i$, the number of kept holes
above $c$, and write $|\psi(c)|=\prod_{p\in Y}|p-c|/p$.

\emph{Factors below $c$.}  $Y\cap[1,c)$ consists of $\{1,\ldots,\ell\}$ and a
set $E\subseteq\{\ell+2,\ldots,c-1\}$ that avoids the $i-2$ holes
$c_2,\ldots,c_{i-1}$, so $|E|=a:=c-\ell-i$.  The factor $(c-p)/p$ is
positive and decreasing in $p$ on $p<c$, so the product over $E$ is at
most the product over the $a$ smallest candidates $\ell+2,\ldots,\ell+1+a$.
Together with $p\le\ell$, and as $\ell+1+a=c-i+1$, this is
$\prod_{p=1}^{c-i+1}(c-p)/p=\binom{c-1}{c-i+1}$, divided by the factor
$(c-\ell-1)/(\ell+1)$ of the missing $p=\ell+1$; as $\binom{c-1}{c-i+1}=\binom{c-1}{i-2}$, this gives
\[
  \prod_{p\in Y,\,p<c}\frac{c-p}p\le\lambda_i(c):=
  \frac{\ell+1}{c-\ell-1}\binom{c-1}{i-2}\quad(i\ge2),\qquad\lambda_1(\ell+1)=1,
\]
the second because for $i=1$ the product is
$\prod_{p=1}^{\ell}(\ell+1-p)/p=1$.

\emph{Factors above $c$.}  $Y\cap(c,N]$ is $(c,N]$ with the $i'$ holes
$c_{i+1},\ldots,c_{\ell+1}$ removed.  The factor $(p-c)/p$ lies in $(0,1)$ and
increases with $p$, so the product is at most the product over the
$N-c-i'$ largest integers of $(c,N]$, which is
$\prod_{p=c+i'+1}^{N}(p-c)/p=\binom{c+i'}c\big/\binom Nc$.  Since
$m\ge c+i'$, we have $N\ge c+i'+q$, and hence
\[
  2^{-c}|\psi(c)|\le U_i(c):=2^{-c}\lambda_i(c)
        \binom{c+i'}c\Big/\binom{c+i'+q}c .
\]

\emph{Moving $c$ down.}  For $i\ge2$ and $c\ge \ell+i$,
\[
  \frac{U_i(c+1)}{U_i(c)}
  =\frac{c-\ell-1}{2(c-\ell)}\cdot\frac c{c-i+2}\cdot\frac{c+i'+1}{c+i'+q+1}
  \le\frac c{2(c-i+2)}\le1,
\]
the last step because $c\ge \ell+i\ge2i-1$ (as $i\le \ell+1$).  Since
$c_i\ge \ell+i$, we get $2^{-c_i}|\psi(c_i)|\le U_i(\ell+i)$ for every $i$; for $i=1$
this is the bound itself.

\emph{The block.}  For $C^\circ=\{\ell+1,\ldots,2\ell+1\}$ with the same $q$, all
three bounds are equalities: $E=\varnothing$, the holes above
$c_i^\circ=\ell+i$ are $\ell+i+1,\ldots,2\ell+1$, and
$N^\circ=2\ell+1+q=c_i^\circ+i'+q$.
Put $Y^\circ=\{1,\ldots,N^\circ\}\setminus C^\circ$ and $\psi^\circ=r_{Y^\circ}$.
Then, term by term,
$2^{-c_i}|\psi(c_i)|\le U_i(\ell+i)=2^{-c_i^\circ}|\psi^\circ(c_i^\circ)|$, so
$\Phi_{\mathrm{hole}}(C,q)\le \Phi_{\mathrm{hole}}(C^\circ,q)$.  By the proof of Lemma~\ref{lem:holeint},
$\Phi_{\mathrm{hole}}(C^\circ,q)=\int_0^1v^q\sum_{c\in C^\circ}|\pi_c|2^{-c}v^{2\ell+1-c}(1-v)^{c-1}\,dv$
with $|\pi_{\ell+i}|=(2\ell+1)!/\bigl(\ell!\,(i-1)!\,(\ell+1-i)!\bigr)$, and the binomial
theorem gives
\[
  \sum_{i=1}^{\ell+1}|\pi_{\ell+i}|2^{-\ell-i}v^{\ell+1-i}(1-v)^{\ell+i-1}
  =\frac{(2\ell+1)!}{\ell!^2\,2^{\ell+1}}(1-v)^\ell\Bigl(v+\frac{1-v}2\Bigr)^\ell
  =\frac{\gamma_\ell}2(1-v^2)^\ell,
\]
where $\gamma_\ell=(2\ell+1)!/(\ell!^24^\ell)$ satisfies
$\int_0^1\gamma_\ell(1-v^2)^\ell\,dv=1$.  The function $(1-v^2)^\ell$ is
nonincreasing and $v^q$ is nondecreasing, so Chebyshev's integral inequality
gives $\Phi_{\mathrm{hole}}(C^\circ,q)\le\frac12\int_0^1v^q\,dv=\frac1{2(q+1)}$.
\end{proof}

\begin{lemma}[The tail beyond the kept holes]\label{lem:holetail}
Let $C,\ell,m,q,N,\psi$ be as in Lemma~\ref{lem:holevals} with $\ell\ge1$
and $q\ge1$.  Then $\Phi_{\mathrm{tail}}(C,q)=\sum_{s>N}|\psi(s)|\,2^{-s}\le\dfrac{\ell+1}{e\,\ell\,(q+1)}$.
\end{lemma}

\begin{proof}
Let $s>N$.  Every $c\in C$ is below $s$, and $c/(s-c)$ increases with $c$.
Since $c_1=\ell+1$ and $c_i\le m-(\ell+1-i)$ for $i\ge2$,
\[
  \prod_{c\in C}\frac c{s-c}
  \le\frac{\ell+1}{s-\ell-1}\prod_{i=2}^{\ell+1}\frac{m-\ell-1+i}{s-m+\ell+1-i}
  =\frac{\ell+1}{s-\ell-1}\binom m\ell\Big/\binom{s-m-1+\ell}\ell .
\]
Apply $1/\binom\nu \ell=(\nu+1)\int_0^1t^\ell(1-t)^{\nu-\ell}\,dt$ with $\nu=s-m-1+\ell$, and
note that $\nu+1=s-m+\ell\le s-\ell-1$ because $m\ge2\ell+1$:
\[
  \prod_{c\in C}\frac c{s-c}\le(\ell+1)\binom m\ell\int_0^1t^\ell(1-t)^{s-m-1}\,dt .
\]
Multiply by $\binom{s-1}N2^{-s}$ and sum over $s>N$.  With
$\sum_{s>N}\binom{s-1}N y^s=\bigl(y/(1-y)\bigr)^{N+1}$ at $y=(1-t)/2$, the
sum of $\binom{s-1}N2^{-s}(1-t)^{s-m-1}$ is
$(1-t)^{-m-1}\bigl((1-t)/(1+t)\bigr)^{N+1}$, so
\[
  \Phi_{\mathrm{tail}}(C,q)\le(\ell+1)\binom m\ell\int_0^1t^\ell
     \Bigl(\frac{1-t}{1+t}\Bigr)^q(1+t)^{-m-1}\,dt .
\]
Now $(1-t)/(1+t)\le e^{-2t}$, since $\log\frac{1+t}{1-t}\ge2t$, and for
$q\ge1$, $(q+1)e^{-2qt}\le2q\,e^{-2qt}\le1/(et)$, since $ye^{-y}\le1/e$.
Extending the integral to $(0,\infty)$,
\[
  (q+1)\,\Phi_{\mathrm{tail}}(C,q)\le\frac{\ell+1}e\binom m\ell\int_0^\infty t^{\ell-1}(1+t)^{-m-1}\,dt
  =\frac{\ell+1}e\binom m\ell\,\mathrm B(\ell,m-\ell+1)=\frac{\ell+1}{e\,\ell}.
\]
\end{proof}

\begin{proposition}[Rows at the top three positions]\label{prop:rowstop}
Let $n\ge1$ and let $A\subset\N_{>0}$ be finite with $\{1,2,3\}\subseteq A$.
Then some finite $Z\supseteq A$ has $|Z|\le|A|+n-1$ and $\Phi(Z)\le1/n$.
\end{proposition}

\begin{proof}
For $n=1$ take $Z=A$ and use Lemma~\ref{lem:confdel}.  Let $n\ge2$, let
$\ell\ge3$ be the leading run of $A$, and let $c_1=\ell+1<c_2<\cdots$ be the
positive integers outside $A$.  Keep $C=\{c_1,\ldots,c_{\ell+1}\}$, fill the
next $n-1$ holes, and put
\[
  N=c_{\ell+n},\qquad Y=\{1,\ldots,N\}\setminus C,\qquad Z=Y\cup A .
\]
Then $Y\supseteq A\cap[1,N]$, the elements of $Z\setminus A$ are the $n-1$
filled holes $c_{\ell+2},\ldots,c_{\ell+n}$, and $|Z|\le|A|+n-1$.  As in the proof
of Proposition~\ref{prop:rowsfree}, adding the elements of $Z$ above $N$ to
$Y$ one at a time never produces a set $\{1,\ldots,t\}$, because the hole
$\ell+1$ lies below them, so Lemma~\ref{lem:confdel} gives
$\Phi(Z)\le\Phi(Y)$.  Put $m=c_{\ell+1}$ and $q=N-m\ge n-1\ge1$.  By
Lemmas~\ref{lem:holevals} and~\ref{lem:holetail},
\[
  \Phi(Y)=\Phi_{\mathrm{hole}}(C,q)+\Phi_{\mathrm{tail}}(C,q)\le\frac1{q+1}\Bigl(\frac12+\frac{\ell+1}{e\,\ell}\Bigr)
  \le\frac1{q+1}\Bigl(\frac12+\frac4{3e}\Bigr)<\frac1{q+1}\le\frac1n,
\]
because $(\ell+1)/\ell\le4/3$ for $\ell\ge3$ and $8<3e$.
\end{proof}

\begin{theorem}[Every row at the root $2$]\label{thm:everyrow}
Let $L\ge1$ and let $F$ be a monic multiple of $(x-2)^L$.  Then
\[
  b_k(F)\ \ge\ L-k+1\qquad(1\le k\le L).
\]
\end{theorem}

\begin{proof}
Since $\deg F\ge L\ge k$, there are at least $k$ nonleading positions.  Let
$S$ be the set of backward distances of $k-1$ positions carrying the $k-1$
largest nonleading magnitudes, and $n=L-k+1$.  Proposition~\ref{prop:rowsfree}
if $\{1,2,3\}\not\subseteq S$, and Proposition~\ref{prop:rowstop} otherwise,
give $Z\supseteq S$ with $|Z|\le|S|+n-1=L-1$ and $\Phi(Z)\le1/n$, and
Lemma~\ref{lem:rowcert} gives $b_k(F)\ge n$.
\end{proof}

The row $k=L$ is $b_L\ge1$, already the row $u=L-1$ of \eqref{eq:order}, and summing the rows gives only
$\Lambda(F)\ge\log L!$, which Proposition~\ref{prop:quadratic} already
exceeds for large $L$; the content of Theorem~\ref{thm:everyrow} is the
individual rows.  The companion paper~\cite{coefficient-mass-rows}
sharpens it to $b_k(F)\ge\beta_2(L-k+1)$, where $\beta_2(n)$ is the least
$b_1$ over monic multiples of $(x-2)^n$, by extending
Lemmas~\ref{lem:rowcert} and~\ref{lem:confdel} to other roots; as
$\beta_2(n)$ is not known in closed form, Theorem~\ref{thm:everyrow} remains
the explicit bound.  For $\ell\in\{1,2\}$ the tail bound of Lemma~\ref{lem:holetail}
exceeds $\frac12$ and the argument of Proposition~\ref{prop:rowstop} does
not close; that is why Proposition~\ref{prop:rowsfree} keeps different
holes there.

\paragraph{Data availability.}
No datasets were generated or analyzed in this study.

\paragraph{Computer assistance.}
No proof in this paper relies on a computation.  The script
\texttt{tests/sweep.py} of the repository, run as a test suite by
\texttt{tests/test\_sweep.py}, checks the following in exact rational
arithmetic with Python's \texttt{fractions} module: the tails
$\frac12$, $\frac3{10}$, $\frac18$ and $\frac9{64}$ of the examples after
Theorem~\ref{thm:tail} and the cutoff example $\Tail(w)=\frac{52}{405}$; the
formula of Lemma~\ref{lem:consecutive} at four rational nodes; that each
$F_n$ of Proposition~\ref{prop:sharp23}, $2\le n\le29$, vanishes at $2$ and
$3$ and has $b_2(F_n)=t_n>2$; and, by an exact linear program, that the least
$b_2$ over monic multiples of $(x-2)(x-3)$ of degree at most $8$ is $t_7$.
It also computes the exact tails of seeded random certificates with two to
four nodes and zeros in $\{1,\ldots,8\}$: in $150$ draws with nodes $1/\rho$,
$\rho\in\{2,\ldots,11\}$, \eqref{eq:tail} and the deletion step in the proof
of Theorem~\ref{thm:tail} always hold, and in $300$ draws that allow nodes
above $\frac12$ both fail for some draws.  The comparison
of~\eqref{eq:jensen} with Theorem~\ref{thm:everyrow} at $L=20$ and $L=40$ in
Section~\ref{sec:bottom} is a floating-point evaluation.  These computations
are checks, not steps of any proof.

\paragraph{Use of AI tools.}
Large language model assistants (Claude, Anthropic) were used in drafting and revising the text, the proofs, the verification code and the Lean formalization. The author has checked the mathematics and is responsible for the content.

\paragraph{Funding and competing interests.}
The author received no specific funding for this work and has no competing
interests to declare.

\begin{thebibliography}{99}
\small
\bibitem{bbbp-multiple}
F.~Beaucoup, P.~Borwein, D.~W. Boyd and C.~Pinner,
\newblock Multiple roots of $[-1,1]$ power series,
\newblock \emph{Journal of the London Mathematical Society} (2) 57 (1998), 135--147,
\newblock \href{https://doi.org/10.1112/S0024610798005857}
{doi:10.1112/S0024610798005857}.

\bibitem{bbbp-ray}
F.~Beaucoup, P.~Borwein, D.~W. Boyd and C.~Pinner,
\newblock Power series with restricted coefficients and a root on a given ray,
\newblock \emph{Mathematics of Computation} 67 (1998), 715--736,
\newblock \href{https://doi.org/10.1090/S0025-5718-98-00960-0}
{doi:10.1090/S0025-5718-98-00960-0}.

\bibitem{bek}
P.~Borwein, T.~Erd\'elyi and G.~K\'os,
\newblock Littlewood-type problems on $[0,1]$,
\newblock \emph{Proceedings of the London Mathematical Society} (3) 79
(1999), 22--46.

\bibitem{bombieri-vaaler}
E.~Bombieri and J.~D. Vaaler,
\newblock Polynomials with low height and prescribed vanishing,
\newblock in A.~C. Adolphson, J.~B. Conrey, A.~Ghosh, and R.~I. Yager (eds.),
\emph{Analytic Number Theory and Diophantine Problems},
Progress in Mathematics 70, Birkh\"auser, 1987, 53--73,
\newblock \href{https://doi.org/10.1007/978-1-4612-4816-3_3}
{doi:10.1007/978-1-4612-4816-3\_3}.

\bibitem{dubickas}
A.~Dubickas,
\newblock Arithmetical properties of powers of algebraic numbers,
\newblock \emph{Bulletin of the London Mathematical Society} 38 (2006), 70--80,
\newblock \href{https://doi.org/10.1112/S0024609305017728}
{doi:10.1112/S0024609305017728}.

\bibitem{dj}
A.~Dubickas and J.~Jankauskas,
\newblock On the reduced height of a polynomial,
\newblock \emph{Publicationes Mathematicae Debrecen} 71 (2007), 325--348,
\newblock \href{https://doi.org/10.5486/PMD.2007.3728}
{doi:10.5486/PMD.2007.3728}.

\bibitem{frw}
M.~Filaseta, M.~L. Robinson and F.~S. Wheeler,
\newblock The minimal Euclidean norm of an algebraic number is effectively
computable,
\newblock \emph{Journal of Algorithms} 16 (1994), 309--333,
\newblock \href{https://doi.org/10.1006/jagm.1994.1015}
{doi:10.1006/jagm.1994.1015}.

\bibitem{grt}
M.~Giesbrecht, D.~S. Roche and H.~Tilak,
\newblock Computing sparse multiples of polynomials,
\newblock \emph{Algorithmica} 64 (2012), 454--480,
\newblock \href{https://doi.org/10.1007/s00453-012-9652-4}
{doi:10.1007/s00453-012-9652-4}.

\bibitem{karlin}
S.~Karlin,
\newblock \emph{Total Positivity, Volume I},
\newblock Stanford University Press, 1968.

\bibitem{mignotte}
M.~Mignotte,
\newblock Sur les multiples des polyn\^omes irr\'eductibles,
\newblock \emph{Bulletin de la Soci\'et\'e Math\'ematique de Belgique} 27
(1975), 225--229.

\bibitem{coefficient-mass-complex}
B.~Pham,
\newblock Coefficient mass of polynomial multiples at complex roots,
\newblock preprint, 2026, \url{https://github.com/bangyen/coefficient-mass}.

\bibitem{coefficient-mass-rows}
B.~Pham,
\newblock Every row is a top row: coefficient order statistics of multiples
of a power,
\newblock preprint, 2026, \url{https://github.com/bangyen/coefficient-mass}.

\bibitem{coefficient-mass-attainment}
B.~Pham,
\newblock Sharpness of the coefficient order-statistic bound,
\newblock preprint, 2026, \url{https://github.com/bangyen/coefficient-mass}.

\bibitem{polya-szego}
G.~P\'olya and G.~Szeg\H{o},
\newblock \emph{Problems and Theorems in Analysis II},
\newblock Classics in Mathematics, Springer, 1998 (reprint of the 1976
  edition), Part Five, Chapter 1,
\newblock \href{https://doi.org/10.1007/978-3-642-61905-2}
{doi:10.1007/978-3-642-61905-2}.

\bibitem{rosser-schoenfeld}
J.~B. Rosser and L.~Schoenfeld,
\newblock Approximate formulas for some functions of prime numbers,
\newblock \emph{Illinois Journal of Mathematics} 6 (1962), 64--94,
\newblock \href{https://doi.org/10.1215/ijm/1255631807}
{doi:10.1215/ijm/1255631807}.

\bibitem{schinzel}
A.~Schinzel,
\newblock On the reduced length of a polynomial with real coefficients,
\newblock \emph{Functiones et Approximatio Commentarii Mathematici} 35
(2006), 271--306,
\newblock \href{https://doi.org/10.7169/facm/1229442629}
{doi:10.7169/facm/1229442629}.

\bibitem{schinzel-ii}
A.~Schinzel,
\newblock On the reduced length of a polynomial with real coefficients, II,
\newblock \emph{Functiones et Approximatio Commentarii Mathematici} 37
(2007), 445--459,
\newblock \href{https://doi.org/10.7169/facm/1229619664}
{doi:10.7169/facm/1229619664}.

\bibitem{schinzel-complex}
A.~Schinzel,
\newblock The reduced length of a polynomial with complex coefficients,
\newblock \emph{Acta Arithmetica} 133 (2008), 73--81,
\newblock \href{https://doi.org/10.4064/aa133-1-5}
{doi:10.4064/aa133-1-5}.
\end{thebibliography}

\end{document}
